\chapter{Validazione sperimentale e benchmark}
\label{chap:validazione-benchmark}

Questo capitolo copre la validazione sperimentale del prototipo consolidato nel capitolo~\ref{chap:progettazione-prototipo}, a partire dalla costituzione di un benchmark di pagine web su cui condurre i test, fino all'analisi dei risultati ottenuti, dei falsi positivi e falsi negativi, e delle \emph{hallucination} prodotte dai modelli linguistici impiegati.

\section{Metodologia di raccolta dei dati e verifica manuale}
\label{sec:metodologia-validazione}

\subsection{Selezione dei casi di studio e raccolta dei dati automatici}
\label{sec:selezione-siti-strumenti}

Il benchmark è composto da 20 pagine web, suddivise in due campioni di pari numerosità.

Il primo campione, di \textbf{10 pagine della \acr{pa}}, è stato selezionato dal patrimonio dati (\emph{open data}) del monitoraggio sull'accessibilità di AGID\footcite{agid:opendata-monitoraggio}, scegliendo pagine appartenenti alle stesse categorie di soggetti monitorati da AGID: comuni, scuole/università e altri enti pubblici. Il secondo campione, di \textbf{10 pagine non-PA}, è stato selezionato fra i domini censiti dal WebAIM Million tramite lo strumento Site Lookup\footcite{webaim:sitelookup}, scegliendo pagine distribuite su categorie eterogenee (\emph{e-commerce}, \emph{news}/\emph{media}, tecnologia, istruzione non-PA\glsadd{pa}, sanità non-PA\glsadd{pa}) in modo da non limitare il confronto a un singolo settore.

Su entrambi i campioni sono stati eseguiti controlli con gli stessi due strumenti automatici usati nei tre \emph{report} di riferimento del capitolo~\ref{chap:metodologia-selezione}, per garantire la comparabilità metodologica dei dati raccolti nel benchmark con quelli dei \emph{report}: \textbf{WAVE\glsadd{wave}}, lo stesso strumento con cui è condotto il WebAIM Million\footcite{wave:extension}, e \textbf{MAUVE\glsadd{mauve}}, lo stesso strumento con cui AGID conduce il monitoraggio semplificato\footcite{mauve:extension}. Entrambi sono stati utilizzati nella forma di estensione del browser, in modo da operare sul DOM effettivamente renderizzato dalla pagina e non su un semplice recupero lato server del codice sorgente.

Nel corso della raccolta è emersa una problematica metodologica ricorrente: su alcune pagine, in particolare quelle con contenuto iniettato lato client o con DOM di dimensioni elevate, l'estensione di MAUVE\glsadd{mauve} sottostima sensibilmente il numero di errori rilevati rispetto a WAVE, restituendo un numero anomalo di criteri \emph{Not evaluated} a fronte di problemi che WAVE individua correttamente sulla stessa pagina. Per questi casi è stato adottato un protocollo di \emph{fallback}: il DOM renderizzato viene copiato manualmente dal pannello \emph{Elements} degli strumenti di sviluppo del browser (F12) e incollato nella modalità di validazione del codice del validatore web ufficiale di MAUVE\glsadd{mauve}\footcite{mauve:single-validation}.

\subsection{Verifica manuale di riferimento}
\label{sec:verifica-manuale-riferimento}

Per stabilire un \emph{ground truth} su cui misurare l'accuratezza del prototipo, ciascuna pagina del benchmark è stata sottoposta, oltre ai controlli automatici di WAVE e \acrg{mauve}, anche a una verifica manuale mirata agli otto criteri \emph{core} e ai tre criteri opzionali individuati nel capitolo~\ref{chap:metodologia-selezione}. La combinazione di controllo automatico e verifica manuale segue l'impostazione della \acr{wcag-em} del W3C, che individua nella verifica umana un passaggio necessario per i casi in cui uno strumento automatico non può accertare la conformità con certezza\footcite{w3c:wcag-em}.

Per ciascun criterio, il protocollo di verifica manuale adottato è derivato direttamente dalle tecniche sufficienti indicate nella documentazione ufficiale del W3C\footcite{w3c:understanding-intro,w3c:quickref}, già richiamate nell'analisi del capitolo~\ref{chap:metodologia-selezione}, qui declinate in un controllo operativo eseguibile su ciascuna pagina del benchmark, con le relative eccezioni e i casi limite più ricorrenti.

\paragraph*{1.1.1 Non-text Content}
\begin{description}
\item[Controllo manuale di riferimento] Per ogni immagine, verificare che il testo alternativo dichiarato (quando presente) trasmetta lo stesso scopo funzionale dell'immagine nel contesto della pagina, distinguendo i casi di immagine decorativa, informativa o funzionale.
\item[Tecniche] \href{https://www.w3.org/WAI/WCAG21/Techniques/general/G94.html}{G94}; \href{https://www.w3.org/WAI/WCAG21/Techniques/html/H37.html}{H37}; \href{https://www.w3.org/WAI/WCAG21/Techniques/aria/ARIA6.html}{ARIA6}.
\item[Note / eccezioni] \texttt{alt=""} è corretto per immagini puramente decorative, non un errore: il fallimento riguarda solo le immagini informative o funzionali prive di alternativa adeguata. Includere anche le immagini di sfondo CSS informative e le immagini-link (l'alt descrive la destinazione, non l'aspetto).
\end{description}

\paragraph*{1.4.3 Contrast (Minimum)}
\begin{description}
\item[Controllo manuale di riferimento] Individuare la coppia testo/sfondo effettivamente resa a schermo quando lo sfondo non è un colore solido dichiarato (immagini, gradienti, stili calcolati a runtime), e campionare il colore realmente visualizzato per applicare la formula di \gls{luminanza-relativa}.
\item[Tecniche] \href{https://www.w3.org/WAI/WCAG21/Techniques/general/G18.html}{G18}; \href{https://www.w3.org/WAI/WCAG21/Techniques/general/G145.html}{G145}.
\item[Note / eccezioni] Esclusi dal criterio: testo di componenti disabilitati, testo puramente decorativo o incidentale, loghi e nomi di marca.
\end{description}

\paragraph*{1.3.1 Info and Relationships}
\begin{description}
\item[Controllo manuale di riferimento] Verificare che intestazioni, liste, tabelle e relazioni label/campo percepibili visivamente siano rappresentate correttamente anche a livello semantico nel markup, non solo in modo sintatticamente valido.
\item[Tecniche] \href{https://www.w3.org/WAI/WCAG21/Techniques/general/G115.html}{G115}; \href{https://www.w3.org/WAI/WCAG21/Techniques/html/H42.html}{H42}; \href{https://www.w3.org/WAI/WCAG21/Techniques/html/H44.html}{H44}; \href{https://www.w3.org/WAI/WCAG21/Techniques/html/H48.html}{H48}; \href{https://www.w3.org/WAI/WCAG21/Techniques/html/H51.html}{H51}.
\item[Note / eccezioni] Verificare anche che l'ordine di lettura programmatico (DOM) corrisponda a quello visivo, quando quest'ultimo è alterato da posizionamento CSS.
\end{description}

\paragraph*{4.1.2 Name, Role, Value}
\begin{description}
\item[Controllo manuale di riferimento] Per ogni componente interattivo non nativo, eseguire l'interazione e verificare che nome accessibile, ruolo e stato (ad esempio \texttt{aria-expanded}, \texttt{aria-checked}) siano esposti e si aggiornino correttamente, non solo che siano dichiarati nel markup statico.
\item[Tecniche] \href{https://www.w3.org/WAI/WCAG21/Techniques/general/G108.html}{G108}; \href{https://www.w3.org/WAI/WCAG21/Techniques/html/H91.html}{H91}; \href{https://www.w3.org/WAI/WCAG21/Techniques/aria/ARIA16.html}{ARIA16}.
\item[Note / eccezioni] Dare priorità ai componenti personalizzati più comuni: menu a tendina, finestre di dialogo/modali, tab, accordion, slider, caroselli.
\end{description}

\paragraph*{2.4.4 Link Purpose (In Context)}
\begin{description}
\item[Controllo manuale di riferimento] Leggere il testo del link nel contesto immediato (frase, paragrafo, cella) e giudicare se lo scopo resta comprensibile anche isolato, come nella modalità \enquote{elenco link} di uno \emph{screen reader}.
\item[Tecniche] \href{https://www.w3.org/WAI/WCAG21/Techniques/general/G91.html}{G91}; \href{https://www.w3.org/WAI/WCAG21/Techniques/html/H30.html}{H30}.
\item[Note / eccezioni] Prestare attenzione ai link generici ripetuti (\enquote{clicca qui}, \enquote{leggi di più}) con testo identico ma destinazioni diverse nella stessa pagina.
\end{description}

\paragraph*{3.3.2 Labels or Instructions}
\begin{description}
\item[Controllo manuale di riferimento] Per ogni campo di un modulo, verificare presenza e corretta associazione dell'etichetta (\texttt{<label for>}, \texttt{aria-label}/\\ \texttt{aria-labelledby}) e che le istruzioni (formato richiesto, campi obbligatori) siano esposte in modo programmatico e non solo visivo.
\item[Tecniche] \href{https://www.w3.org/WAI/WCAG21/Techniques/general/G131.html}{G131}; \href{https://www.w3.org/WAI/WCAG21/Techniques/html/H90.html}{H90}.
\item[Note / eccezioni] Pattern di fallimento più frequente: \texttt{placeholder} usato al posto di un'etichetta persistente (scompare alla digitazione) — è un fallimento anche se il campo resta comprensibile a campo vuoto.
\end{description}

\paragraph*{2.4.7 Focus Visible}
\begin{description}
\item[Controllo manuale di riferimento] Navigare l'intera pagina esclusivamente da tastiera (\emph{Tab}/\emph{Shift+Tab}) su ogni elemento interattivo e verificare visivamente che l'indicatore di focus sia presente e sufficientemente distinguibile dallo stato non attivo.
\item[Tecniche] \href{https://www.w3.org/WAI/WCAG21/Techniques/general/G149.html}{G149}; \href{https://www.w3.org/WAI/WCAG21/Techniques/css/C15.html}{C15}.
\item[Note / eccezioni] Controllare anche gli elementi focusabili che un altro elemento sovrapposto (ad esempio per \emph{z-index}) può occultare, pur essendo l'indicatore tecnicamente presente nel markup.
\end{description}

\paragraph*{2.4.1 Bypass Blocks}
\begin{description}
\item[Controllo manuale di riferimento] Verificare non solo la presenza tecnica di un meccanismo di salto (link \enquote{skip to content}, \emph{landmark} ARIA/HTML5), ma la sua reale utilizzabilità: visibilità al focus, correttezza della destinazione, presenza coerente su tutte le pagine del sito.
\item[Tecniche] \href{https://www.w3.org/WAI/WCAG21/Techniques/general/G1.html}{G1}; \href{https://www.w3.org/WAI/WCAG21/Techniques/aria/ARIA11.html}{ARIA11}.
\item[Note / eccezioni] La coerenza su più pagine è verificabile solo se il benchmark include più pagine per sito; sulla pagina campionata si verificano comunque presenza e usabilità del meccanismo.
\end{description}

\paragraph*{1.4.1 Use of Color}
\begin{description}
\item[Controllo manuale di riferimento] Individuare i casi in cui un'informazione (stato di errore, campo obbligatorio, link nel testo) è veicolata solo dal colore, senza un secondo indicatore non cromatico (icona, sottolineatura, testo).
\item[Tecniche] \href{https://www.w3.org/WAI/WCAG21/Techniques/general/G14.html}{G14}; \href{https://www.w3.org/WAI/WCAG21/Techniques/general/G182.html}{G182}.
\item[Note / eccezioni] Include anche l'uso del solo colore per distinguere serie di dati in grafici o voci di legenda.
\end{description}

\paragraph*{1.4.10 Reflow}
\begin{description}
\item[Controllo manuale di riferimento] Ridimensionare il \emph{viewport} a 320px (equivalente a zoom 400\%) e verificare l'assenza di scorrimento bidirezionale e di perdita di contenuto o funzionalità.
\item[Tecniche] \href{https://www.w3.org/WAI/WCAG21/Techniques/css/C32.html}{C32}; \href{https://www.w3.org/WAI/WCAG21/Techniques/general/G206.html}{G206}.
\item[Note / eccezioni] Sono esentati dal reflow i contenuti che richiedono necessariamente una disposizione bidimensionale per essere compresi (tabelle di dati, mappe, diagrammi, editor con barra degli strumenti).
\end{description}

\paragraph*{1.4.11 Non-text Contrast}
\begin{description}
\item[Controllo manuale di riferimento] Stessa logica di 1.4.3, applicata però a componenti d'interfaccia ed elementi grafici essenziali (bordi di campo, icone informative, indicatore di focus).
\item[Tecniche] \href{https://www.w3.org/WAI/WCAG21/Techniques/general/G195.html}{G195}; \href{https://www.w3.org/WAI/WCAG21/Techniques/general/G207.html}{G207}.
\item[Note / eccezioni] Stesse esclusioni di 1.4.3 per componenti disabilitati o puramente decorativi.
\end{description}

Per la maggior parte dei criteri il controllo manuale di riferimento è stato eseguito da chi scrive, ispezionando pagina e markup nel browser; per i criteri 2.4.7 e 1.4.10, che richiedono l'osservazione di uno stato dinamico dell'interfaccia (rispettivamente: navigazione da tastiera e ridimensionamento del \emph{viewport}), la verifica è stata condotta interagendo direttamente con la pagina resa, non sul solo codice sorgente statico.


\section{Casi di studio e composizione del benchmark}
\label{sec:casi-studio-benchmark}

\todoplaceholder{elenco e descrizione dei casi di studio che compongono il benchmark, con i criteri di definizione dei singoli casi (non le cifre aggregate, già riportate in \S\ref{sec:selezione-siti-strumenti}).}

% TODO: come indicato nel Piano di lavoro, la definizione dei casi di studio
% non è stata vincolata a un singolo progetto di riferimento. In particolare,
% l'ipotesi iniziale di utilizzare il progetto universitario "Student Space"
% come caso di studio è stata scartata su indicazione della Prof.ssa Gaggi, in
% quanto Student Space presenta un livello di accessibilità già
% sufficientemente buono e conterrebbe quindi pochi problemi reali da
% individuare; l'introduzione artificiale di errori non è stata considerata
% una soluzione adeguata, poiché errori sintetici potrebbero non rappresentare
% correttamente quelli riscontrabili nei siti reali. Il benchmark è stato
% quindi costituito prendendo a riferimento: (1) le principali problematiche
% di accessibilità evidenziate dal report WebAIM Million, già discusso nel
% capitolo~\ref{chap:metodologia-selezione}; (2) la letteratura di settore;
% (3) problematiche rappresentative riscontrabili in pagine web reali,
% relative agli otto criteri \emph{core} selezionati.
%
% NB: numerosità (20 pagine, 10 PA + 10 non-PA), fonti di selezione (AGID
% opendata, WebAIM Site Lookup) e categorie di sito sono già descritte in
% \S\ref{sec:selezione-siti-strumenti} (non ripetere qui). Questa sezione
% deve riportare solo l'elenco concreto dei singoli siti/pagine selezionati,
% con indicazione per ciascuno del criterio o dei criteri \emph{core}
% rappresentati, della fonte puntuale (voce specifica del catalogo AGID o
% dominio del WebAIM Million) e dell'eventuale motivazione di inclusione
% (ad esempio un problema noto in letteratura che il sito esemplifica).


\section{Strategia ed esecuzione dei test}
\label{sec:strategia-test}

\todoplaceholder{esecuzione del prototipo/LLM sul benchmark e metriche adottate per il confronto con il \emph{ground truth}.}

% TODO: la raccolta automatica (WAVE/MAUVE, \S\ref{sec:selezione-siti-strumenti})
% e il protocollo di verifica manuale (\S\ref{sec:verifica-manuale-riferimento})
% sono già descritti: non ripeterli qui. Questa sezione deve coprire solo ciò
% che manca: il protocollo con cui il prototipo/LLM è stato eseguito su
% ciascuna pagina del benchmark (stesso input degli strumenti automatici?
% stesso ambito dei criteri \emph{core}/opzionali?), e le metriche adottate
% (accuratezza, precisione, richiamo, F1 o altre pertinenti) per confrontare
% l'esito del prototipo con il \emph{ground truth} stabilito in
% \S\ref{sec:verifica-manuale-riferimento}.


\section{Analisi dei falsi positivi e falsi negativi}
\label{sec:analisi-fp-fn}

\todoplaceholder{analisi sistematica dei falsi positivi e falsi negativi riscontrati durante la validazione.}

% TODO (settimana 4 del piano di lavoro): analisi sistematica dei casi in cui
% il sistema ha segnalato un problema di accessibilità inesistente (falso
% positivo) o non ha rilevato un problema realmente presente (falso
% negativo), per ciascuno degli otto criteri \emph{core}, con discussione delle cause
% ricorrenti individuate.


\section{Analisi delle hallucination degli LLM}
\label{sec:analisi-hallucination}

\todoplaceholder{analisi delle hallucination prodotte dai modelli linguistici nel contesto della verifica di accessibilità.}

% TODO (settimana 4 del piano di lavoro): analisi dei casi in cui l'LLM ha
% prodotto un'informazione non fondata sul contenuto reale della pagina
% (ad esempio una motivazione plausibile ma tecnicamente errata, o un
% riferimento a un elemento non presente), con discussione dell'impatto di
% questo fenomeno sull'affidabilità della verifica assistita da LLM.


\section{Risultati sperimentali e discussione}
\label{sec:risultati-sperimentali}

\todoplaceholder{sintesi dei risultati quantitativi e qualitativi della validazione, con relativa discussione.}

% TODO: documentazione dei risultati (settimana 4 del piano di lavoro):
% sintesi quantitativa dei risultati per criterio e complessiva; confronto,
% dove pertinente, fra i modelli LLM già messi a confronto nel
% capitolo~\ref{chap:progettazione-prototipo}; discussione dei limiti
% riscontrati e del grado di affidabilità raggiunto dal sistema.